\Needspace{23\baselineskip}\vspace{6mm}
\SourceLeaf{586}\pdfbookmark[1]{Astronomische Nachrichten Nr. 417. 1841. Februar 9.}{nachrichten-417}
\begin{center}\small\textsc{Schumacher}.\quad Astronomische Nachrichten Nr.\ 417.\quad 1841.\quad Februar 9.\end{center}

In einem öffentlichen Blatte fand ich unlängst die Nachricht, dass der Amerikanische Marine-Capitain \textsc{Wilkes}
\begin{center}\textit{dem magnetischen Südpole}\end{center}
ziemlich nahe gekommen sei, und dass er in $67^0\;4'$ südl.\ Breite und $147^0\;30'$ Länge (ohne Zweifel östlich von Greenwich) die magnetische Abweichung $12^0\;35'$ östlich, und die Neigung $87^0\;30'$ gefunden habe. Nach einer flüchtig angestellten Rechnung würde ich nun hienach einstweilen den wirklichen Pol
\begin{center}in $70^0\;21'$ südlicher Breite,\qquad $146^0\;17'$ Länge\end{center}
setzen. Dieser Platz liegt demjenigen, welchen meine Theorie [S.\ 163 d.\ B.] angegeben hat, viel näher, als ich selbst erwartet hatte. Der wirkliche Pol, wie ich dort vermuthet hatte, nordlicher, als der nach der Theorie berechnete; aber der Unterschied in der Breite erreicht nur den dritten Theil von dem, auf welchen ich nach Ansicht der Beobachtungen von Hobarttown gefasst war. Eben so liegt der wirkliche Pol westlicher, als der nach der Theorie berechnete, und hier ist der Unterschied fast genau so gross wie der a.\ a.\ O.\ von mir präsumirte. Übrigens ist unnöthig zu bemerken, dass in diesen hohen Breiten der Unterschied von sechs Längengraden nur eben so viel bedeutet, wie zwei Breitengrade.

\clearpage
\SourceLeaf{587}\pdfbookmark[0]{Anzeigen nicht eigner Schriften.}{anzeigen-nicht-eigner}
\thispagestyle{empty}
\vspace*{35mm}
\begin{center}
{\LARGE ANZEIGEN}\par\vspace{16mm}
{\large NICHT EIGNER}\par\vspace{16mm}
{\LARGE SCHRIFTEN.}
\end{center}
\clearpage
\SourceLeaf{589}\fancyhead[L]{\small ANZEIGEN NICHT EIGNER SCHRIFTEN.}
\pdfbookmark[1]{Ueber die Daltonsche Theorie. 1830.}{daltonsche-theorie}
\begin{center}\small Göttingische gelehrte Anzeigen.\quad 1830 December 11.\end{center}

\textit{Ueber die \textup{\textsc{Dalton}}sche Theorie, von \textup{J. F.\ \textsc{Benzenberg}.} Düsseldorf, \textup{1830.} Bei J. E.\ Schaub.} (192 Seiten in 8., nebst drei Steindrucktafeln).

Die von \textsc{Dalton} aufgestellte Hypothese, dass die verschiedenen Gasarten, aus welchen die atmosphärische Luft besteht, gar nicht gegenseitig auf einander drücken, sondern eben so viele von einander gleichsam unabhängige Atmosphären ausmachen, hat bei wenigen Physikern bisher Beifall gefunden. Unter diesen zeichnet sich aber Herr \textsc{Benzenberg} durch den unermüdeten warmen Eifer, mit welchem er jene Hypothese seit beinahe zwanzig Jahren in Schutz nimmt, ganz besonders aus. Namentlich hat er in der \textsc{Daubuisson}schen trigonometrisch-barometrischen Messung der Höhe des Monte Gregorio einen wichtigen Grund \textit{für} die \textsc{Dalton}sche Hypothese gefunden. Es ist klar, dass die barometrischen Höhenmessungen, wenn die \textsc{Dalton}sche Hypothese wahr ist, anders berechnet werden müssen, als nach der gewöhnlichen Theorie. Bei dem 5260 Fuss hohen Monte Gregorio fand Herr \textsc{Benzenberg} das Resultat der ersten Rechnung um 16 Fuss kleiner, als nach der andern, und sehr nahe eben so viel übertraf letztere das Resultat der trigonometrischen Messung, welche Differenz mithin nach Herrn \textsc{Benzenberg}s Rechnung durch die Annahme der \textsc{Dalton}schen Hypothese fast vollkommen gehoben werden würde. Herr \textsc{Benzenberg} hat diese Rechnungen zuerst in \textsc{Gilbert}'s Annalen der Physik 1812 bekannt gemacht, und ist auch nachher
\SourceLeaf{590}an andern Orten zu wiederholten Malen damit aufgetreten. Auch über andere Abschnitte der Physik, welche mit der \textsc{Dalton}schen Vorstellungsart in Berührung kommen, wie die Akustik und Eudiometrie, hat Herr \textsc{Benzenberg} sich ausgelassen, nicht sowohl, um Gründe \textit{für} jene Hypothese darin zu suchen, als vielmehr, um diejenigen Gründe, welche man daraus gegen dieselbe hernehmen kann und hergenommen hat, zu bekämpfen. In gegenwärtiger Schrift ist alles abermals vereinigt aufgestellt.

Die Anerkennung. welche diese Bemühungen bei den Physikern und Mathematikern gefunden haben, scheint, wenigstens was die ausländischen betrifft, nicht genügend gewesen zu sein. Herr \textsc{Benzenberg} erzählt in der Vorrede seines Buchs, dass er bei seinem Aufenthalt in Paris im Jahre 1815 \textsc{Gilberts} Annalen in der Bibliothek des Instituts gar nicht, und das Exemplar in der grossen Königlichen Bibliothek nicht aufgeschnitten gefunden habe. Mit \textsc{Laplace} habe er von seiner Theorie der Berechnung der Barometerhöhen in \textsc{Daltons} Hypothese gesprochen; allein dieser grosse Geometer sei damals schon alt gewesen.

Den deutschen Physikern kann doch der Vorwurf der Nichtbeachtung nicht mit Recht gemacht werden. Wir lesen in der neuen Ausgabe des physikalischen Wörterbuchs im Artikel Atmosphäre, dass `Herr \textsc{Benzenberg} der bedeutendste, gründlichste und eifrigste Vertheidiger der \textsc{Dalton}schen Theorie ist, dass mit Anwendung derselben die Höhen genauer, als ohne sie, berechnet werden, und dass darin ein bedeutender Beweis für die Richtigkeit derselben liege.'

Allein \textit{geprüft} scheint der Verfasser dieses Artikels die \textsc{Benzenberg}schen Rechnungen nicht zu haben: alle Zahlen sind nur ohne weiteres aus den \textsc{Gilbert}schen Annalen copirt. Dasselbe gilt von demjenigen, was über jene Rechnungen in dem Artikel Höhenmessung in dem erwähnten Wörterbuche gesagt ist. Vielleicht haben die Verfasser beider Artikel eine Prüfung deswegen für minder wesentlich gehalten, weil sie, den barometrischen Höhenmessungen überhaupt eine geringere Zuverlässigkeit beilegend, als Herr \textsc{Benzenberg}, die Beweiskraft von dessen Rechnungen doch nicht anerkannten, obwohl freilich der Verfasser des ersten Artikels dadurch das vorher angeführte zum Theil wieder aufhebt.

Es kann nicht die Absicht der gegenwärtigen Anzeige sein, unsere eigene Ansicht von der Zulässigkeit oder Unzulässigkeit der \textsc{Dalton}schen Hypothese selbst zu entwickeln, sondern wir beschränken uns auf dasjenige, was Herr \textsc{Ben\SourceLeaf{591}zenberg} zur Unterstützung derselben in vorliegender Schrift von neuem vorgetragen hat, und namentlich auf seine Berechnung der barometrischen Höhenmessungen.

Es schien vor allen Dingen nöthig, erst die Richtigkeit dieser \textsc{Benzenberg}schen Berechnung selbst zu prüfen. Zu unserer Verwunderung ist daraus hervorgegangen, dass diese Berechnung unrichtig ist, und dass eine richtig geführte Rechnung ein ganz entgegengesetztes Resultat gibt.

Es wird nöthig sein, dies hier mit einiger Ausführlichkeit und so weit nachzuweisen, dass jeder in den Stand gesetzt wird, selbst zu urtheilen, um so mehr, da es sich hier nicht sowohl um einen Rechnungsfehler im eigentlichen Sinn, als um einen Irrthum im Räsonnement handelt.

Auf der zweiten und dritten Seite des Buchs finden wir zwei tabellarische Übersichten, jede von drei Columnen, welche wir einzeln durchgehen müssen (dieselben stehen in \textsc{Gilberts} Annalen der Physik, B.\ 42.\ S.\ 163 und 164).

In der ersten Columne wird, nach \textsc{Dalton}, die in 100 Theilen trockner atmosphärischer Luft enthaltene
\begin{center}
\begin{tabular}{lclrl}
Stickluft & .\quad .\quad . & zu & $78.93$ & Theilen\\
Sauerstoffluft & . & — & $21.00$ & —\\
Kohlensaure Luft & & — & $0.07$ & —
\end{tabular}
\end{center}
alles dem Raume nach, angegeben (Statt der zweiten Zahl steht im Buche selbst $21.90$, welches ein offenbarer Druckfehler ist).

Die zweite Columne setzt die specifischen Gewichte
\begin{center}
\begin{tabular}{lcr}
der Stickluft & .\quad .\quad . & $0.9691$\\
der Sauerstoffluft & .\quad . & $1.1148$\\
der kohlensauren Luft & & $1.5000$
\end{tabular}
\end{center}
das der gemeinen atmosphärischen Luft zur Einheit angenommen. Die erste und dritte Zahl sind von \textsc{Biot} entlehnt; die zweite hat Herr \textsc{Benzenberg}, wie er selbst sagt, so berechnet, `dass die Mischung genau das von \textsc{Biot} gegebene Gewicht trockner atmosphärischer Luft habe, nemlich $\frac{1}{10493}$ des Quecksilbers bei $0^0$ Reaumur und 28 Zoll Druck.' Man sieht, dass Herr \textsc{Benzenberg} sich selbst nicht klar gemacht hat, was er hier eigentlich hat rechnen wollen, denn das eben angezeigte ist \textit{dieser} Rechnung fremd, und offenbar kam es blos darauf an, der Mi%
